\documentclass[11pt,a4paper]{article}

\usepackage[brazil]{babel}
\usepackage[utf8]{inputenc}
\usepackage[T1]{fontenc}
\usepackage{lmodern}
\usepackage{geometry}
\usepackage{setspace}
\usepackage{parskip}
\usepackage{titlesec}
\usepackage{xcolor}
\usepackage{hyperref}

\geometry{margin=2.5cm}
\setstretch{1.2}

\definecolor{primary}{RGB}{41,128,185}
\hypersetup{
    colorlinks=true,
    linkcolor=primary,
    urlcolor=primary
}

\titleformat{\section}{\Large\bfseries\color{primary}}{\thesection}{1em}{}
\titleformat{\subsection}{\large\bfseries}{\thesubsection}{1em}{}

\title{\textbf{Plano de Gestão \\ GULECD/UFPR}}
\author{Chapa 01 - "stats.sh"}
\date{\today}

\begin{document}

\maketitle

\section{Introdução}

O presente Plano de Gestão tem como objetivo delinear as diretrizes estratégicas, operacionais e institucionais da chapa Chapa 01 - "stats.sh" para a coordenação do Grupo de Usuários Linux e Software Livre do curso de Estatística e Ciência de Dados da Universidade Federal do Paraná (GULECD/UFPR).

Este plano está fundamentado nos princípios estabelecidos pelo Manifesto, pela Carta de Fundação e pelo Código de Conduta do grupo, buscando consolidar o GULECD/UFPR como um ambiente de excelência técnica, colaboração aberta e desenvolvimento coletivo.

\section{Visão Geral da Gestão}

A proposta desta chapa está orientada para a consolidação institucional do grupo, expansão de suas atividades e fortalecimento de sua comunidade interna e externa. Pretende-se estabelecer uma base sólida que permita a continuidade e expansão das iniciativas do GULECD/UFPR.

A visão da chapa é edificar um grupo com bases sólidas e duradouras, uma agremiação autossustentável que sobreviva e cresça além da passagem de cada indivíduo, construindo um legado coletivo.

\section{Objetivos Estratégicos}

\subsection{Institucionalização e Reconhecimento}

\begin{itemize}
    \item Formalizar o GULECD/UFPR como projeto de extensão oficial da UFPR de forma que suas iniciativas e ações sejam reconhecidas como horas formativas.
    \item Obter reconhecimento institucional junto a departamentos, coordenações de curso e demais órgãos relevantes da universidade.
    \item Estabelecer o grupo como referência em software livre/código aberto e estatística/ciência de dados no âmbito institucional e nacional.
    \item Formar parceria com um \textit{"host fiscal"} visando capacidade de receber auxílios, patrocínios e demais formas de recursos destinados a atividades do grupo.
\end{itemize}

\subsection{Estruturação do Ambiente de Desenvolvimento}

\begin{itemize}
    \item Estruturar um ambiente institucional que permita aos membros desenvolver projetos de forma independente e colaborativa.
    \item Implementar o \textbf{LD-ID (Laboratório Digital de Iniciativas Discentes)}, voltado ao desenvolvimento de software, análises estatísticas, aprendizado de máquina e big data.
    \item Garantir infraestrutura adequada para suporte a projetos acadêmicos e experimentais.
    \item Buscar recursos institucionais e parcerias para viabilizar a abertura contínua dos laboratórios A e B do PA para uso dos alunos.
\end{itemize}


\subsection{Parcerias e Sustentabilidade}

\begin{itemize}
    \item Estabelecer parcerias com empresas de tecnologia visando benefícios aos membros, tais como:
    \begin{itemize}
        \item créditos em serviços de computação em nuvem;
        \item acesso a plataformas de desenvolvimento e ferramentas profissionais;
        \item programas de capacitação técnica e certificação;
        \item apoio a eventos e iniciativas do grupo.
    \end{itemize}
\end{itemize}

\subsection{Integração com a Comunidade de Software Livre}

\begin{itemize}
    \item Inserir o GULECD/UFPR na comunidade de software livre em níveis local, nacional e internacional.
    \item Estabelecer conexões com outros grupos, comunidades e organizações de software livre.
\end{itemize}

\section{Atividades e Iniciativas}

\subsection{Formação e Capacitação}

\begin{itemize}
    \item Promover cursos e palestras sobre uso e desenvolvimento de software livre e de código aberto.
    \item Desenvolver material educativo estruturado e acessível.
    \item Criar trilhas de formação técnica contínua.
    \item Implementar um sistema de mentoria entre membros.
\end{itemize}

\subsection{Eventos e Comunidade}

\begin{itemize}
    \item Organizar \textit{install parties} para incentivo à adoção de Linux.
    \item Promover encontros presenciais periódicos em ambiente descontraído, com o objetivo de fortalecer a integração social e o senso de comunidade entre os membros.
    \item Desenvolver a \textbf{1ª Conferência de Software Livre e Código Aberto em Estatística e Ciência de Dados}.
    \item Organização de hackathons e sprints de desenvolvimento.
    \item Realizar a "Pesquisa Anual de Uso de Tecnologias Digitais" com o objetivo de inferir o perfil de uso de tecnologias digitais por parte dos alunos do Curso de Estatística e Ciência de Dados da UFPR.
\end{itemize}

\subsection{Projetos Estruturantes}

\begin{itemize}
    \item Expandir o projeto da Biblioteca de Distribuições Linux.
    \item Criar e manter canais de suporte ao uso de Linux e software livre.
    \item Desenvolver projetos colaborativos de software e análise de dados.
    \item Desenvolver website oficial do grupo.
    \item Desenvolver website voltado para compartilhamento de notícias em software livre/código aberto e estatística/ciência de dados
\end{itemize}

\section{Governança e Operação}

\begin{itemize}
    \item Manter estrutura organizacional semi-horizontal.
    \item Garantir transparência em todas as decisões e atividades.
    \item Incentivar a participação ativa dos membros nas decisões do grupo.
\end{itemize}

\section{Considerações Finais}

Este Plano de Gestão visa tanto a manutenção das atividades do GULECD/UFPR quanto a sua expansão e institucionalização, com sustentabilidade, relevância acadêmica e impacto na formação dos estudantes.

A chapa compromete-se a atuar com responsabilidade, transparência e alinhamento aos princípios do grupo, promovendo um ambiente colaborativo e tecnicamente sólido.

\end{document}